\makeatletter
% Fonts declarations
% Use first for different initials
\newfontfamily{\initials}[Extension=.ttf,ExternalLocation=fonts/,Scale=MatchUppercase]{Kanzlei-Initialen}
\newfontfamily{\sectionFont}[Extension=.ttf,ExternalLocation=fonts/]{Germanica}
%\newfontfamily{\numerals}[Extension=.ttf]{roman_font_7.ttf}
\newfontfamily{\titleFont}[Extension=.ttf,ExternalLocation=fonts/]{MediciText}

% Table of contents format. An example of how it can be changed for different parts. 

%\titlecontents{chapter}
%[0em]                 % adjust left margin
%{\chapterFont\LARGE}             % font formatting
%{} % section label and offset
%{}
%{\titlerule*[0.4em]{.}\normalfont\contentspage}

%\titlecontents{section}
%[1em]                 % adjust left margin
%{\articleFont\large}             % font formatting
%{} % section label and offset
%{}
%{\titlerule*[0.4em]{.}\normalfont\contentspage}

%\titlecontents{subsection}
%[2em]                 % adjust left margin
%{\bfseries}             % font formatting
%{\thecontentslabel\hspace{1em}} % section label and offset
%{}
%{\titlerule*[0.4em]{.}\normalfont\contentspage}

%\titlecontents{subsubsection}
%[3em]                 % adjust left margin
%{}             % font formatting
%{\thecontentslabel\hspace{1em}} % section label and offset
%{}
%{\titlerule*[0.4em]{.}\normalfont\contentspage}

%\renewcommand\tableofcontents{%
%    \begin{multicols}{2}[\section*{\contentsname
%        \@mkboth{%
%           \MakeUppercase\contentsname}{\MakeUppercase\contentsname}}]%
%    \@starttoc{toc}%
%    \end{multicols}%
%    }
%\renewcommand\contentsname{Table of contents}
% \setcounter{tocdepth}{3} 

% Sections formatting
\titleformat{\chapter}[display]
    {\vfill\filcenter\bf}
    {}
    {0ex}
    {\rput[b](0pt,20pt){\includegraphics[width=9cm]{formatting/ornament}}\\
    \fontsize{29}{23}\sectionFont#1\\
    \rput[b]{180}(0pt,-5pt){\includegraphics[width=9cm]{formatting/ornament}}}
    [\vfill\null\thispagestyle{empty}\clearpage]
\titlespacing{\chapter}{0pt}{0ex}{0ex}

\titleformat{\section}[block]
            {\centering\Large}
            {\centering\MakeUppercase{\bf\roman{section}}\\}
            {0pt}
            {\sectionFont#1}
\renewcommand*\l@section{\@dottedtocline{1}{1.5em}{3.3em}}

% Fancy frame to be used
% Page background decoration
\newlength{\outerMargin}
\newlength{\innerMargin}
% Change these 3 values
% A6: 0.3cm, 0.8cm
\setlength{\outerMargin}{0.8cm}
\setlength{\innerMargin}{1.3cm}
\newcommand\scaleFactor{1}

\newlength{\frameSpaceA}
\setlength{\frameSpaceA}{0.132cm}
\newlength{\frameSpaceB}
\setlength{\frameSpaceB}{0.1985cm}

\newlength{\cornerOffsetBase}
\setlength{\cornerOffsetBase}{0.2145pt}
\newlength{\cornerOffsetImage}
\setlength{\cornerOffsetImage}{0.078cm}

\definecolor{frameGrey}{RGB}{128,128,128}


\backgroundsetup{
    color=black,
    scale=1,
    opacity=1,
    angle=0,
    contents={
        \begin{tikzpicture}
        \fill [fill=white] (current page.south east) rectangle (current page.north west)(current page.north west);
        \ifodd\value{page}
            % External border
        \draw[line width=4pt*\scaleFactor]
            ( $ (current page.south east) + (-\outerMargin,\outerMargin) $ ) rectangle ( $ (current page.north west) + (\innerMargin,-\outerMargin) $ );
	\draw[line width=1.2pt*\scaleFactor]
            ( $ (current page.south east) + (-\outerMargin - \frameSpaceA * \scaleFactor, \outerMargin + \frameSpaceA * \scaleFactor) $ ) rectangle ( $ (current page.north west) + (\innerMargin + \frameSpaceA * \scaleFactor, -\outerMargin - \frameSpaceA * \scaleFactor) $ );
	\draw[line width=0.8pt*\scaleFactor]
            ( $ (current page.south east) + (-\outerMargin - \frameSpaceB * \scaleFactor, \outerMargin + \frameSpaceB * \scaleFactor) $ ) rectangle ( $ (current page.north west) + (\innerMargin + \frameSpaceB * \scaleFactor, -\outerMargin - \frameSpaceB * \scaleFactor) $ );
            % Corner ornaments
            \node[anchor=north east, rotate=90] at ( $ (current page.north west) + (\innerMargin,-\outerMargin ) + (-\cornerOffsetImage*\scaleFactor-\cornerOffsetBase,\cornerOffsetImage*\scaleFactor+\cornerOffsetBase)$ ){\includegraphics[width=1.5cm*\scaleFactor]{formatting/corner6_b}};
            \node[anchor=north east, rotate=0] at ( $ (current page.north east) + (-\outerMargin, -\outerMargin) + (\cornerOffsetImage*\scaleFactor+\cornerOffsetBase,\cornerOffsetImage*\scaleFactor+\cornerOffsetBase) $ ){\includegraphics[width=1.5cm*\scaleFactor]{formatting/corner6_b}};
            \node[anchor=north east, rotate=180] at ( $ (current page.south west) + (\innerMargin,\outerMargin) + (-\cornerOffsetImage*\scaleFactor-\cornerOffsetBase,-\cornerOffsetImage*\scaleFactor-\cornerOffsetBase) $ ){\includegraphics[,width=1.5cm*\scaleFactor]{formatting/corner6_b}};
            \node[anchor=north east, rotate=270] at ( $ (current page.south east) + (-\outerMargin,\outerMargin) + (\cornerOffsetImage*\scaleFactor+\cornerOffsetBase, -\cornerOffsetImage*\scaleFactor-\cornerOffsetBase) $ ){\includegraphics[width=1.5cm*\scaleFactor]{formatting/corner6_b}};
        \else
            	\draw[line width=4pt*\scaleFactor]
            ( $ (current page.south east) + (-\innerMargin, \outerMargin) $ ) rectangle ( $ (current page.north west) + (\outerMargin, -\outerMargin) $ );
	  	\draw[line width=1.2pt*\scaleFactor]
            ( $ (current page.south east) + (-\innerMargin - \frameSpaceA * \scaleFactor, \outerMargin + \frameSpaceA * \scaleFactor) $ ) rectangle ( $ (current page.north west) + (\outerMargin + \frameSpaceA * \scaleFactor, -\outerMargin - \frameSpaceA * \scaleFactor) $ );
		\draw[line width=0.8pt*\scaleFactor]
            ( $ (current page.south east) + (-\innerMargin - \frameSpaceB * \scaleFactor, \outerMargin + \frameSpaceB* \scaleFactor) $ ) rectangle ( $ (current page.north west) + (\outerMargin + \frameSpaceB * \scaleFactor, -\outerMargin - \frameSpaceB * \scaleFactor) $ );
            % Corner ornaments
            \node[anchor=north east, rotate=90] at ( $ (current page.north west) + (\outerMargin,-\outerMargin ) + (-\cornerOffsetImage*\scaleFactor-\cornerOffsetBase,\cornerOffsetImage*\scaleFactor+\cornerOffsetBase)$ ){\includegraphics[width=1.5cm*\scaleFactor]{formatting/corner6_b}};
            \node[anchor=north east, rotate=0] at ( $ (current page.north east) + (-\innerMargin, -\outerMargin) + (\cornerOffsetImage*\scaleFactor+\cornerOffsetBase,\cornerOffsetImage*\scaleFactor+\cornerOffsetBase) $ ){\includegraphics[width=1.5cm*\scaleFactor]{formatting/corner6_b}};
            \node[anchor=north east, rotate=180] at ( $ (current page.south west) + (\outerMargin,\outerMargin) + (-\cornerOffsetImage*\scaleFactor-\cornerOffsetBase,-\cornerOffsetImage*\scaleFactor-\cornerOffsetBase) $ ){\includegraphics[width=1.5cm*\scaleFactor]{formatting/corner6_b}};
            \node[anchor=north east, rotate=270] at ( $ (current page.south east) + (-\innerMargin,\outerMargin) + (\cornerOffsetImage*\scaleFactor+\cornerOffsetBase, -\cornerOffsetImage*\scaleFactor-\cornerOffsetBase) $ ){\includegraphics[width=1.5cm*\scaleFactor]{formatting/corner6_b}};
        \fi
        \end{tikzpicture}
    }
}

% Optional line between columns, with ornament included
%\SetMCRule{
%    custom-line={
%        \path ($ (TOP) !.4! (BOT) $) to [ornament=88] ($ (TOP) !.6! (BOT) $);
%    },
%    width=0.2pt,
%    extend-fill
%}

% Footer formatting
\fancyhf{}
\renewcommand{\headrulewidth}{0pt}
\setlength{\footskip}{15pt}
\cfoot{{\rput[b]{90}(-0.2,0.15){\pgfornament[width=0.5cm]{26}} \thepage \rput[b]{-90}(0.2, 0.03){\pgfornament[width=0.5cm]{26}}}}

% I have no idea what this was for
\newenvironment{chapquote}[2][2em]
{
	\def\chapquote@author{#2}%
	\itshape
	``}
{''\par\normalfont\hfill--\ \chapquote@author\hspace*{\@tempdima}\par\bigskip}

% Main font. Linlibertine includes a good amount of UTF-8 characters e.g. greek ones
\setmainfont[ExternalLocation=fonts/]{linlibertine}[
    SizeFeatures={Size=13},
	Extension=.ttf,
	UprightFont=*-re,
	ItalicFont=*-it,
	BoldFont=*-bd,
	BoldItalicFont=*-bi.ttf
]

% Big capital font
\renewcommand{\LettrineFontHook}{\initials}
% Use after a paragraph with a big capital letter if the next paragraph overlaps with the capital letter
\newcommand{\checklettrine}{
	\ifnum\prevgraf<\lettrinelines
        \wlog{\prevgraf}
		\Repeat{\fpeval{\lettrinelines-\prevgraf}}{\hfill \\}
		\indent
	\fi
}
\setcounter{DefaultLines}{\lettrinelines}
\renewcommand*{\DefaultLoversize}{0.1}
\makeatother

% Paragraph 
\setlength{\parskip}{0.6em}
\renewcommand{\baselinestretch}{1.3}

% Use these if you want to manually update or reset section/subsection/whatever counters for the table of contents. In this example subsections are only reset by chapters, which do not reset sections
%\newcommand{\mychapter}[2]{
%    \chapter*{#2}
%    \setcounter{subsection}{0}
%    \addcontentsline{toc}{chapter}{#1}
%}
%
%\newcommand{\mysection}[1]{
%    \section*{#1}
%    \addcontentsline{toc}{section}{#1}
%}
%\counterwithout{subsection}{section}